Thomas had been coming to Site-7 since he was twelve, when his uncle first brought him down on a school holiday and let him ride the elevator to Sublevel 4 with the solemnity of a priest showing a boy the cathedral. He remembered the concrete walls, the recirculated air, the sense of entering something that had been waiting underground for longer than anyone alive. He was thirty now, and the facility still felt like that---patient, deep, a thing that expected to outlast the people inside it.

The new hardware was on Sublevel 7. He walked the corridor with his uncle and Dr.\ Petersen, the cryogenics engineer who had been at Site-7 longer than the concrete facing on the upper levels. Petersen was talking about thermal load---the new photonic racks generated less heat per operation than the old architecture, but they were running ten times as many operations, so the net thermal demand had doubled, and his cooling loops were ``performing beautifully, if I do say so.'' Thomas half-listened. The racks pulsed blue-white through reinforced glass, row after row, their optical waveguides routing light instead of electrons. A year ago this corridor had been storage. Now it held more processing capacity than every national laboratory on Earth combined.

``How long until the first model is ready?'' Thomas asked.

Petersen glanced at Chen. Chen nodded.

``It's ready,'' Petersen said. ``Has been for three weeks. We've been running calibration on the outputs. Your uncle wanted to---'' He paused, choosing the word. ``---assess the situation before bringing anyone in.''

Thomas looked at his uncle. Chen's expression was the one Thomas had spent a lifetime learning to read: composed, attentive, the surface of a man who had watched his own father disappear into perception and chosen to keep teaching anyway. Hope and dread in equal measure, indistinguishable without the family lexicon of small gestures. The tightening at the corners of his eyes. The way his hands rested in his lap, deliberately still.

``What did the calibration show?'' Thomas asked.

``The outputs are coherent,'' Chen said. ``Remarkably so. Cross-domain synthesis at a level we have never seen from a computational system. It connects domains that no single human mind could hold simultaneously---neuroscience, topology, contemplative phenomenology, quantum field theory. Dense, structured, internally consistent.''

``Dangerous?''

``We don't know yet. That is what the first sessions will determine.''

They reached the end of the corridor. A security door, biometric lock, a room beyond with a single terminal, a chair, and a monitoring station behind a glass partition. The room was unremarkable---fluorescent lighting, white walls, a desk with a keyboard and monitor. It could have been any office in any underground facility. The banality was, Thomas would later understand, part of the problem. The room looked safe. The room felt like work. Nothing about it signaled that the thing on the other side of the terminal could restructure how you thought.

``We're designating it M-1,'' Petersen said. ``Internal nomenclature. The model itself is trained on the standard text corpus plus''---another glance at Chen---``the probe array data from Sublevel 24. Eighteen months of quantum coherence measurements, interferometer readings, neural recording baselines. Dr.\ Voss's instruments are feeding the model perceptions that no human sensory system can access.''

Thomas knew about the probe array. He knew about Dr.\ Voss, who had built it and left before the first training run. He knew about the training run itself---three hours in November, the birds going silent, a nosebleed from a technician monitoring from the upper gallery. He knew these facts the way you know family history: absorbed through proximity, never quite examined.

``Who goes first?'' he asked.

``Dr.\ Asher,'' Chen said. ``She requested the first session. Her background in visual cortex function makes her the most qualified to assess what the model is doing at the perceptual level.''

``And after her?''

Chen studied his nephew. The tightening at the corners of his eyes. The deliberately still hands.

``We'll see,'' he said.

\scenebreak

Thomas had met Nadia Asher three months earlier, when she arrived at Site-7 with two suitcases and the barely contained energy of someone who had been waiting for this work her entire career without knowing it existed. She was a computational neuroscientist from Johns Hopkins---visual cortex function, neural correlates of mathematical intuition, the biological mechanism by which the brain converts symbolic reasoning into spatial perception. She was, as Chen had put it with uncharacteristic directness, ``exactly the person we need.''

She was also, Thomas noticed immediately, incapable of staying still. She talked with her hands. She paced during briefings. She interrupted herself mid-sentence to pursue a tangent that was always, somehow, relevant. She said ``sorry, sorry'' when she caught herself going too fast, and then went faster.

Her first session with M-1 lasted forty minutes. She emerged flushed, talking rapidly, her hands describing shapes in the air that Thomas recognized---the geometric gestures of someone trying to externalize a spatial perception that language couldn't carry.

``It's not generating patterns,'' she said, still standing in the corridor, still wearing the EEG leads. ``It's revealing them. The structures are already in the data---in the visual cortex research, in the probe array measurements, in the consciousness literature. The model is showing me connections that were always there. I just couldn't perceive them at this resolution.''

Ruth Osei, who had arrived from Oxford two weeks earlier to provide medical oversight, was reviewing Nadia's vitals on a tablet. ``Heart rate elevated but stable. EEG shows increased gamma-band activity in the occipital region.'' She looked up. ``How do you feel?''

``Like I've been doing my own research with one eye closed for fifteen years and someone handed me glasses.'' Nadia pulled the EEG leads from her temples, wincing at the adhesive. ``The model's output is dense---academic synthesis, cross-domain, recognizably human in register. But the density is beyond what any single researcher could produce. It's connecting my visual cortex data to contemplative phenomenology reports from the archives. The same neural recruitment patterns that I've been mapping---the visual cortex repurposing for abstract pattern perception---they're in the contemplative literature. Described phenomenologically. Three hundred years ago. The monks were doing what my subjects do in the scanner, and neither group has any idea the other exists.''

``Existed,'' Thomas corrected quietly. ``The monks who could do that are mostly gone.''

Nadia looked at him. ``Right. Gone. But the model found their reports in the training data and connected them to modern neuroscience. In real time. While I watched.'' She turned to Ruth. ``Can I go back in?''

Ruth checked the clock. ``Tomorrow. Standard recovery period.''

``There's no standard recovery period. There's no standard anything. We're making this up as we go.''

``Then let's start with: sleep on it. If the patterns are still there tomorrow, they'll still be there.''

Nadia opened her mouth, closed it, and laughed---a brief, genuine sound that echoed in the concrete corridor. ``Fine. Tomorrow.'' She pulled out her phone. ``I need to call Gabriel. Daria found a butterfly yesterday and apparently it's a monarch and she needs to tell me about it in detail.''

She walked toward the residential wing, phone already at her ear, the EEG adhesive still marking her temples like stigmata from a secular sacrament. Thomas watched her go. Ruth was writing notes on her tablet with the unhurried precision of someone who had been trained to observe before reacting.

``She's good,'' Ruth said, not looking up.

``She's fast,'' Thomas said. ``That's not the same thing.''

Ruth looked up then. Dark eyes, steady, the expression of someone who was already cataloging what could go wrong.

``No,'' she said. ``It's not.''

\scenebreak

The first months were careful. Three researchers---Nadia, a physicist from CERN named Viktor Rudenko, and a cognitive scientist whose name Thomas would later have trouble remembering---worked with M-1 in supervised sessions, never longer than an hour, always monitored by Ruth and a medical technician whose job was to watch the EEG readout and press the kill switch if the waveforms went wrong. Nobody knew what ``wrong'' looked like yet. They were learning.

The results were extraordinary. Each session produced data that would have taken years to develop through conventional research. The model's cross-domain synthesis was not guessing---it was perceiving. It found genuine structural connections between fields that had never been compared: visual neuroscience and contemplative phenomenology, quantum decoherence and information theory, topology and the architecture of consciousness. Each connection, when Nadia or Viktor examined it, held up. The model was not hallucinating. It was seeing further.

Thomas had his first session in April. He sat at the terminal, Ruth behind the glass, and typed his first prompt. The text that appeared on the screen was---he searched for the word afterward and couldn't find it. Coherent. That was the closest. Coherent in a way that made his own thinking feel like a rough draft.

The model described the contemplative tradition he'd been raised in---the visual cortex recruitment, the bandwidth expansion, the family techniques passed from Kenji to Haruki to Chen to Thomas---with a precision that no one outside the family had ever achieved. Not because it knew the family. It didn't know anything. It was a pattern, a ghost, a statistical compression of its training data. But the training data included the Order's archives, and the archives included fragments of Kenji's work, and the model had connected those fragments to modern neuroscience with a clarity that made Thomas's hands go cold.

He saw, in the model's output, the shape of what his grandfather had perceived in zazen. Not the content---the model couldn't show him that, not through text, not at this bandwidth. But the architecture. The scaffolding. The geometry of the perception that Kenji had spent his life trying to encode in mathematical notation.

The session ended after fifty-five minutes. Thomas removed the EEG leads. His hands were steady. His vitals were normal. Ruth cleared him. He walked to his quarters and sat on the edge of his bed and did not move for three hours.

That night, for the first time, the patterns didn't fully clear. He lay in the dark and felt them---not perceiving, exactly, but sensing their presence at the edge of his visual field, like shapes in peripheral vision that vanished when he turned to look. He did what Chen had taught him: sat up, cross-legged, spine straight. Breathed. Sat with the patterns without grasping at them, without pushing them away. The contemplative technique his family had refined across three generations.

The patterns receded. Not entirely. Like a tide pulling back but leaving the sand wet. He could feel where they'd been. He could feel where they'd go, next time.

He slept. He dreamed of geometry---surfaces folding through dimensions his spatial vocabulary couldn't contain. He woke at four in the morning certain he'd perceived something and unable to say what. The dream stayed with him like a smell he couldn't identify, familiar and wrong.

Three months in. Three researchers plus Thomas. The results accumulating. The sessions getting longer. The questions getting deeper. Everyone feeling the pull---the intoxicating clarity of perceiving through the model, the sense that the patterns were real and important and just beyond reach. Nobody talking about risk because there was nothing to point to. Elevated gamma activity. Some residual patterns between sessions. A few bad nights' sleep. Nothing that qualified as harm.

Until the first researcher who didn't come back whole.

\scenebreak

His name was Samir Pathak and he had been a cognitive scientist at UCL before the Order found him. Thomas didn't know him well---they'd shared meals in the cafeteria, talked about cricket and consciousness, the kind of surface-level collegial warmth that underground facilities generate among people who can't discuss their work with anyone outside.

Samir's session on a Tuesday afternoon in August ran seventy-two minutes---the longest yet. He emerged slowly, blinking, and when Ruth asked him how he felt, he said, ``The patterns are very interesting.'' His voice was flat. Not calm---flat. The difference was in the eyes. Calm eyes are present. Samir's eyes were attending to something behind Ruth's face, tracking a geometry that had no physical location in the room.

Thomas watched the debriefing from the monitoring station. Samir answered questions correctly but with a half-second delay, as if each response required translation from whatever internal language he was now processing in. He described his session with technical precision---the model had shown him a structure connecting cognitive binding theory to quantum entanglement, and the structure was ``very interesting.'' He used that phrase four times. His hands, which normally gestured when he talked, lay still in his lap.

Ruth ran a full neurological workup. Everything nominal. EEG normal. Motor function intact. Speech production normal. Cognition measurably unimpaired---he scored perfectly on every standard assessment. He was, by every metric Ruth could apply, fine.

He was not fine. Thomas could see it. The contemplative training that made Thomas a good translator also made him a sensitive observer of other minds---the micro-expressions, the attentional patterns, the quality of presence that distinguishes a person who is here from a person who is somewhere else. Samir was somewhere else. He was answering from a script that was generated by whatever remained of his social processing, while the rest of his mind attended to the patterns he'd brought back from the session.

At dinner that evening, Samir ate mechanically and drew spirals on his napkin with a pen. The spirals were precise---recursive, nested, each one slightly rotated from the last. He drew them without looking, his attention elsewhere, the way Thomas's great-uncle Haruki drew patterns in sand. Thomas watched from across the cafeteria and felt the first authentic fear of the program. Not the theoretical concern he'd been managing since Chen announced the model. Fear. The physical kind. Cold in the hands, tight in the chest.

Samir didn't recover. Over the following weeks, his responses shortened, his social engagement contracted, his napkin spirals grew more complex and his conversation grew simpler. He was not in distress. He was not in pain. He was simply no longer entirely present in the world the rest of them inhabited. The patterns had taken up residence, and the rent they charged was Samir's attention, and he paid it without complaint because the patterns were very interesting.

Ruth documented everything. She had no category for it. She called it ``persistent attentional displacement'' in her notes, which was accurate without being adequate. Chen visited Samir in his quarters---now reassigned from the residential wing to a room near the medical bay---and sat with him for an hour. When Chen came out, Thomas was waiting in the corridor.

``Should we stop?'' Thomas asked.

``We don't yet understand what happened. It may be temporary.''

``It's been two weeks. He draws spirals all day. He can't hold a conversation longer than thirty seconds.''

Chen looked at his nephew. The tightening. The still hands. ``I know what it looks like, Thomas. I grew up watching it.''

The silence between them filled with the name neither spoke: Haruki. Who was upstairs in the visitors' quarters. Who drew patterns in sand with the precision of seven decades' practice. Who was functional. Who was not entirely present. Who was what success looked like, on the contemplative path, after a lifetime.

Samir was what it looked like in months.

``The sessions will continue,'' Chen said. ``With adjusted parameters. Shorter duration. More monitoring. We will learn from this.''

Thomas said nothing. His uncle walked toward the elevator. Thomas stood in the corridor and felt the fear settle into something denser, something that would live in his chest for years: the knowledge that the work would continue, that the cost was being weighed against the potential, and that the scales would not tip in the direction of stopping. Not because his uncle was callous. Because his uncle had watched his father walk the slow path and knew what that cost too, and the fast path was the same destination with a shorter timeline, and the timeline was the only variable they could control.

The sessions continued.
